% 
% ###################################################################
%  Copyright (c) 2025, AUTHORNAME 
%  Editors: A.D. Rollett & M. De Graef
%  All rights reserved.
%
% Licensed under the Creative Commons CC BY-NC-SA 4.0 License, 
% hereafter referred to as the "License"; you may not use this 
% document except in compliance with the License. You may obtain 
% a copy of the License at 
%     https://creativecommons.org/licenses/by-nc-sa/4.0/legalcode 
% Unless required by applicable law or agreed to in writing, all 
% material distributed under the License is distributed on an 
% "AS IS" BASIS, WITHOUT WARRANTIES OR CONDITIONS OF ANY KIND, 
% either express or implied. See the License for the specific 
% language governing permissions and limitations under the License.
% ###################################################################

% ###################################################################
% ###################################################################
% ###################################################################
% The following lines are to be uncommented or edited as needed 

%     uncomment this line only if this chapter is a core/foundational chapter;
%     for a core chapter a "Foundational" label will appear on the top left above the chapter title.
%\corechapter{Yes}

%     this will appear in a secondary header below the chapter title.
\OCchapterauthor{Marc De Graef, Carnegie Mellon University}

% All figures are stored in the src/ChapterTitle/eps folder and must be of the *.eps, *.pdf, or *.png type. 

% this is the 6-character (all uppercase) chapter descriptor used in labels and refs...
\renewcommand{\chabbr}{MICROD}

%     replace \noheaderimage by the chapter header image file name (without .eps extension).
%     pdf files are also acceptable.
%     Chapter header images must be 2480 x 1240 pixels with 300dpi, RGB format.
\chapterimage{\noheaderimage}

% start the chapter and define the chapter label (outside of the \chapter command!)
\chapter{Microstructure Descriptors}\label{\chabbr:MicrostructureDescriptors} % should be the same as the filename (without extension)

% add the author information so that it will appear in the Author List at the start of the document
% repeat this line for each author
\writeauthor{MICROD:MicrostructureDescriptors}{Microstructure Descriptors}{De Graef}{Marc}{Materials Science \& Engineering}{Carnegie Mellon University}{mdg@andrew.cmu.edu}{https://kernie.materials.cmu.edu/~degraef/www/}

% Each chapter begins with Learning Objectives; the list of objectives should have links to sections/subsections using their OC labels.
% Each Learning Objective should be an active statement (i.e., contain a verb).  The \\ command can be used to force an item into the 
% second column if LaTeX breaks the line at an awkward location.  The label color (OCBurntOrange) should not be changed.
% Within the itemize environment, hyphenation is temporarily turned off.
\begin{learningobjectives}
% This itemized list should be formatted as follows and must be enclosed by { } 
{\begin{itemize}
    \item[{\color{OCBurntOrange}\ref{\chabbr:whatisit}:}] Understand what a microstructure is
%    \item[{\color{OCBurntOrange}\ref{\chabbr:quaternions}:}] Define the quaternion number system\\
%    \item[{\color{OCBurntOrange}\ref{\chabbr:algebras}:}] Learn about Normed Division Algebras
%    \item[{\color{OCBurntOrange}\ref{\chabbr:othernumbers}:}] Discover other number systems, such as dual numbers and split numbers
\end{itemize}}
\end{learningobjectives}
% ###################################################################
% ###################################################################
% ###################################################################


%%%%%%%%%%%%%%%%%%%%%%%%%
%%%%%%%%%%%%%%%%%%%%%%%%%
% This where the actual chapter content starts   %
%%%%%%%%%%%%%%%%%%%%%%%%%
%%%%%%%%%%%%%%%%%%%%%%%%%

\section{Introductory Remarks}\label{\chabbr:sec:whatisit}
Modern materials characterization tools cover the entire range of length scales encountered in materials, as illustrated in Fig.~\ref{\chabbr:fig:lengthscales}; from the atomistic range (using scanning and transmission electron microscopes, x-ray diffraction, and scanning probe microscopes) to the macroscopic range, visible to the unaided eye.  In this chapter, we introduce a number of descriptors that are mostly applicable to the intermediate ``micro'' range which covers length scales between about $100$ nm and $1$ mm.  In this ``micro'' range, we can use visible light (often in the form of an optical microscope) to study what the material looks like. Accordingly, we call the collective features that are ``visible'' at this length scale range the \indexit{microstructure} of the material; in other words, microstructure refers to those features that can be made visible using, essentially, an optical microscope (or an instrument that uses a similar wave length range).   In broad terms, we distinguish between \indexit{single crystalline materials} and \indexit{polycrystalline materials}. In the former case, the crystal lattice has a single orientation across the entire sample, in the latter case the sample consists of many domains with different orientations of the crystal lattice; we call such domains \indexit{grains}. 

In this chapter, we will focus our attention on polycrystalline materials, and we will introduce a series of grain descriptors, ranging from grain size and grain shape to grain orientation.  We will see that all of those grain attributes have an impact of material properties and behavior when the material is subjected to a set of external conditions (temperature, stress, electromagnetic fields, etc.). The analysis of grain orientations leads to the study of a material's \indexit{texture} and this will be the topic of a separate chapter.

Quantification of microstructures is an important area in MSE; we borrow techniques and tools from statistics, image recognition and processing, and mathematics, and we combine those tools into algorithms that are suitable to solve the microstructure problems we face.  Machine learning tools are becoming increasingly relevant to address specific problems, such as feature segmentation, background corrections, and contrast improvement.  But before we can use advanced tools, we must begin with the basic microstructure descriptors that underly many of our theoretical models for microstructure properties: grain size, area/volume fraction, grain boundary length and area, and so on.  Historically, optical microscopy played an important role in this area; 2D images were readily obtainable (after proper sample preparation), and before the advent of computer-based image processing, standardized manual measurement methods were developed along with techniques to connect the measured 2D quantities to their 3D equivalents.  This is the domain of \indexit{stereology}, a branch of mathematics, and we will cover some of the basic stereological concepts in this chapter.  We then conclude this chapter with a brief introduction to 2D and 3D shape quantification.

\begin{figure}[t!]
\centering\leavevmode
\includegraphics[width=6in]{\graphicspath lengthscales.eps}
\caption{\label{\chabbr:fig:lengthscales}Microstructure length scales, the electromagnetic spectrum, and the resolution ranges of various experimental characterization techniques.}
\end{figure}

\section{Grain Size Measurements}\label{\chabbr:sec:grainsize}
In this section, we will consider \indexit{equiaxed grains}, i.e., grains that have very nearly the same size when measured in random directions. The perfect equiaxed 2D shape would obviously be the circle, but it is not possible to tile the plane without gaps using circles, nor is it possible to tile all of space with spheres.  Grains touch each other along \indexit{grain boundaries} and, in general, three boundaries will meet in a \indexit{triple junction}. In the absence of significant porosity, there is in general no empty space between grains. Equiaxed grains can generally be well approximated by a circle (or sphere, in 3D); this means that we can assign an \indexit{equivalent circle diameter} as an approximate grain size.  One of the central questions to be answered is then: how do we measure the diameter of a large number of grains without introducing a statistical bias in the measurements? As we will see in the next subsection, it is not uncommon for such microstructure measurements to be carried out using random sampling.

\subsection{Random Measurements}\label{\chabbr:ssec:random}
The French naturalist and mathematician, the Comte de Buffon (1707-1788), proposed the following problem: consider a plane, ruled with equidistant parallel lines, where the distance between the lines is $D$.  A needle of length $L<D$ is tossed onto the plane. What is the probability that the needle intersects one of the lines?  The situation is illustrated in Fig.~\ref{\chabbr:fig:buffon}.  Needles of length $L$ are randomly tossed onto the grid of lines in part (a) of the figure.  If we consider one particular case, shown in (b), then we see that this needle will intersect the top line, provided the distance $L\sin\theta$ is greater than $y$.  If $L\sin\theta$ is less than $y$, then the needle is not long enough to intersect the top line.  The range of $y$ values is $[0,D]$, and the range of $\theta$ values is $[0,\pi]$, so we can draw a diagram which shows all possible positions and orientations for the needle (the gray rectangle in Fig.~\ref{\chabbr:fig:buffon}c) and all cases for which the needle intersects a line, given by the area below the sine-curve.  The probability of randomly hitting a line is then equal to the area under the sine-curve, divided by the area of the rectangle, or:
\[
	P = \frac{L}{\pi D} \int_0^{\pi} \sin\theta\,\mathrm{d}\theta = \frac{2L}{\pi D}.
\]

\begin{figure}[b!]
\centering\leavevmode
\includegraphics[width=6in]{\graphicspath buffon.eps}
\caption{\label{\chabbr:fig:buffon}(Microstructure length scales, the electromagnetic spectrum, and the resolution ranges of various experimental characterization techniques.)a) schematic of the Buffon needle experiment; (b) geometry of the problem; (c) integration area to obtain the probability $P$.}
\end{figure}

If we now carry out this experiment for a needle with length $L=D$, then the probability is precisely $P=2/\pi$. In practice, we would toss the needle a number of times, say $N$, and count how many times it intersects a line, say $C$, and then we would have $\pi=2N/C$, an estimate of $\pi$!  It is therefore possible to determine the value of $\pi$, by randomly tossing a needle onto a set of parallel lines.\footnote{Over a nine-year period, we had the students of the junior level class on Microstructure and Properties at Carnegie Mellon University carry out $100$ attempts each and count how many times an intersection occurred. The total number of attempts over that period was $N=27,600$ and the number of intersections was $C=17,478$, leading to an estimate for $\pi$ of $3.1583$, i.e. about $0.5\%$ too large.}  Table~\ref{tb:buffon} shows how closely we can approximate $\pi$ by repeating the experiment.  Obviously, these are numerical results, not from a real experiment!\footnote{There are several websites which have Java applets illustrating this problem;  a good one can be found at \href{http://mste.illinois.edu/activity/buffon/}{{\color{blue}here}}.}


\begin{table}[t!]
\centering\leavevmode
\begin{tabular}{|c|c|l|r|}
\hline
$N$ & $C$ &  $2N/C$ & \% error\\
\hline\hline
$10^{2}$ & 63&3.17460& 1.0507\\
$10^{3}$ & 652 &3.06748& -2.3589\\
$10^{4}$ & 6504 & 3.07503& -2.1187\\
$10^{5}$ & 63338 &3.15766& 0.5115\\
$10^{6}$ & 636091 & 3.14420& 0.0831\\
$10^{7}$ & 6364338 &3.14251& 0.0292\\
$10^{8}$ & 63662298 & 3.14158&-0.0005\\
\hline
Exact & $2/\pi=0.63661977$&3.1415926&---\\
\hline
\end{tabular}
\caption{\label{tb:buffon}Results from simulations based on uniform random number generation of the Buffon needle problem.}
\end{table}

% !TEX root = ../main.tex
% ###################################################################
% Copyright (c) 2020, Marc De Graef 
% All rights reserved.
%
% Licensed under the Creative Commons CC BY-NC-SA 4.0 License, 
% hereafter referred to as the "License"; you may not use this 
% document except in compliance with the License. You may obtain 
% a copy of the License at 
%     https://creativecommons.org/licenses/by-nc-sa/4.0/legalcode 
% Unless required by applicable law or agreed to in writing, all 
% material distributed under the License is distributed on an 
% "AS IS" BASIS, WITHOUT WARRANTIES OR CONDITIONS OF ANY KIND, 
% either express or implied. See the License for the specific 
% language governing permissions and limitations under the License.
% ###################################################################

% ###################################################################
% The following lines are to be uncommented or edited as needed 

%\corechapter{Yes}
%     uncomment this line only if this chapter is a core/foundational chapter;
%     for a core chapter a "Core" label will appear on the top left above the chapter title.

\chapterauthor{Marc De Graef, Carnegie Mellon University}
%     this will appear in a secondary header below the chapter title.

% All figures are stored in the src/MaterialProperty/eps folder and must be of the *.eps type. 
\renewcommand{\chaptergraphicspath}{../Figures/}

\chapterimage{Figures/chap5header}
%     replace \noheaderimage by the chapter header image file name (without .eps extension).
%     Chapter header images must be 2480 x 1240 pixels with 300dpi, RGB format.
% ###################################################################



\chapter{Microstructure Descriptors\label{chap:descriptors}}

\begin{messagebox}{What you will learn in this chapter}{ocre}{\icinfo}{white}
\begin{itemize}
	\item Geometrical descriptors for microstructures (\ref{sec:geometry})
	\item Basic principles of stereology (\ref{sec:stereology})
	\item One-point and two-point correlation functions (\ref{sec:correlations})
	\item Aspects of 3D shape analysis (\ref{sec:3Dshapes})
\end{itemize}
\end{messagebox}

\section{Introduction}\label{sec:DescriptorIntroduction}
It is one of the goals of materials science and engineering to predict how a given microstructure will behave when subjected to a set of external conditions (temperature, stress, electromagnetic fields, etc.).  To do so, we must be able to quantify the microstructure, either from 2D images or from 3D reconstructions.   This is a complex problem because we may not always know which quantities are most relevant; is it the grain size that matters most for a material's strength or should we be measuring the total surface area of the grains?  Does the volume fraction of a second phase matter, or is its spatial distribution more important?  What about the shape of precipitates and inclusions? 

Quantification of microstructures is an important area in MSE; we borrow techniques and tools from statistics, image recognition and processing, and mathematics, and we combine those tools into algorithms that are suitable to solve the microstructure problems we face.  Machine learning tools are becoming increasingly relevant to address specific problems.  But before we discuss advanced tools, we must begin with the basic microstructure descriptors that underly many of our theoretical models for microstructure properties: grain size, area/volume fraction, grain boundary length and area, and so on.  Historically, optical microscopy played an important role in this area; 2D images were readily obtainable (after proper sample preparation), and before the advent of computer-based image processing, manual measurement methods were developed along with techniques to connect 2D quantities to their 3D equivalents.  This is the domain of \indexit{stereology}, a branch of mathematics, and we will cover some of the basic stereological concepts in this chapter.  We conclude this chapter with a brief introduction to 3D shape quantification.

\newpage
Modern materials characterization tools span the entire range of length scales encountered in materials, as illustrated in Fig.~\ref{fig:lengthscales}, from the atomistic range (using scanning and transmission electron microscope and scanning probe microscopes) to the macroscopic range.  The descriptors introduced in the following sections are mostly applicable to the ``micro'' range which covers length scales between about $100$ nm and $1$ mm.

\begin{figure}[t]
\centering\leavevmode
\includegraphics[width=4in]{Figures/lengthscales}
\caption{Microstructure length scales, the electromagnetic spectrum, and the resolution ranges of experimental techniques.}
\label{fig:lengthscales}
\end{figure}






\section{Geometrical Microstructure Descriptors}\label{sec:geometry}

\begin{messagebox}{Important Section}{ocre}{\icattention}{black}
In this section, we introduce several geometrical parameters that can be used to quantify a 2D microstructure image.  We also cover a few important software tools that can be used for these computations. In addition, we will show that sometimes random measurements can be used to determine quantitative descriptors.
\end{messagebox}

\subsection{Buffon's needle problem}
In this section, we will consider an \indexit{equiaxed} microstructure, i.e., a microstructure in which the grains have approximately the same dimensions in each direction.  This excludes microstructures with long, elongated grains, which are typically produced when the material is plastically deformed in a particular direction.  The first question we will attempt to answer is: \textit{How do we measure the grain size of a three-dimensional structure, if we only have a two-dimensional image?}  In a typical microstructural observation, we would use an optical microscope with a given magnification $M$, and obtain an image of the microstructure that contains a relatively large number of individual grains.  The surface of the sample is a 2D section through the true microstructure, and it is rather difficult to imagine what the 3D shape of each grain is without actually cutting it out of the material (which is usually not an option).  So we have to look at the grains in the image from a more pragmatic point of view, which can be achieved by means of statistical models.  \indexit{Stereology} provides the tools needed to determine 3D quantities from lower dimensional sections.  We will introduce the topic of stereology, which is covered in more detail in section~\ref{sec:stereology}, by means of the famous \textit{Buffon's Needle Problem}.\index{Buffon's needle problem}

\begin{figure}[t]
\includegraphics[width=5.9in]{Figures/buffon}
\caption{Illustration of the Buffon Needle Problem.}
\label{fig:buffon}
\end{figure}

\indexit{Buffon's needle} is a problem first posed by the French naturalist and mathematician, the Comte de Buffon (1707-1788). Consider a plane, ruled with equidistant parallel lines, where the distance between the lines is $D$.  A needle of length $L<D$ is tossed onto the plane. What is the probability that the needle intersects one of the lines?

The situation is illustrated in Fig.~\ref{fig:buffon}.  Needles of length $L$ are randomly tossed onto the grid of lines (a).  If we consider one particular case (b), then we see that this needle will intersect the top line, provided the distance $L\sin\theta$ is greater than $y$.  If $L\sin\theta$ is less than $y$, then the needle is not long enough to intersect the top line.  The range of $y$ values is $[0,D]$, and the range of $\theta$ values is $[0,\pi]$, so we can draw a diagram which shows all possible positions and orientations for the needle (the gray rectangle in Fig.~\ref{fig:buffon}c) and all cases for which the needle intersects a line (area below the sine-curve).  The probability of randomly hitting a line is then equal to the area under the sine-curve, divided by the area of the rectangle, or:
\[
	P = \frac{2L}{\pi D}.
\]
If we now carry out this experiment for a needle with length $L=D$, then the probability is precisely $P=2/\pi$. In practice, we would toss the needle a number of times, say $N$, and count how many times it intersects a line, say $C$, and then we would have $\pi=2N/C$, an estimate of $\pi$!  It is therefore possible to determine the value of $\pi$, by randomly tossing a needle onto a set of parallel lines.\footnote{This is an easy experiment to carry out yourself.  You can find several websites which have Java applets illustrating this problem;  a good one can be found at \href{http://mste.illinois.edu/activity/buffon/}{here}.}  Table~\ref{tb:buffon} shows how closely we can approximate $\pi$ by repeating the experiment.  Obviously, this is a numerical result, not a real experiment!

\begin{table}[h]
\centering\leavevmode
\begin{tabular}{|c|l|r|}
\hline
\# attempts & $2N/C$ & \% error\\
\hline\hline
$10^{1}$ & 2.22222& -29.264\\
$10^{2}$ & 3.27869& 4.364\\
$10^{3}$ & 3.29489& 4.880\\
$10^{4}$ & 3.14514& 0.113\\
$10^{5}$ & 3.14767& 0.193\\
$10^{6}$ & 3.14229& 0.022\\
$10^{7}$ & 3.14299& 0.044\\
$10^{8}$ & 3.14129&-0.010\\
\hline
\end{tabular}
\caption{\label{tb:buffon}Results from simulations of the Buffon needle problem.}
\end{table}


It is not unusual to use \indexit{random sampling} approaches to determine or estimate quantities that would otherwise be difficult to measure and Buffon's needle problem serves as a particularly simple example of this approach. Random sampling lies at the core of a class of numerical techniques commonly known as \indexit{Monte Carlo simulations}


\subsection{Manual grain size measurements}
In the materials field, several techniques have been introduced for measuring the grain size based on random sampling, mostly before the advent of digital image processing.  While these methods are based on random sampling, they are precisely defined by the \indexit{American Society for Testing of Materials} (ASTM)\index{ASTM}, in their document E112: \textit{Standard test methods for determining average grain size}. There are several methods described in this document, and we will select one of them: the \indexit{lineal intercept method}. Consider a hypothetical grain structure, such as the one shown in Fig.~\ref{fig:grains}.  The magnification of the ``image''  is $100\times$, so that $1$ mm in the image corresponds to $10 \mu$m in the sample.  The lineal intercept method is carried out as follows:






%\subsection{}\label{\chabbr:}

















